Begins with the non-divergent boundary layer form of the compressible equations, 

\begin{equation} 
\frac{\partial \hat{U}}{\partial x}+
\frac{\partial \hat{V}}{\partial y}
=0
,
\end{equation} 

for two  species, the oxidizer (O)  

\begin{equation} 
 \hat{U}\frac{\partial \psi_{o}}{\partial x}+
 \hat{V}\frac{\partial \psi_{o}}{\partial y}
=
\frac{1}{PrLe_o}\frac{\partial}{\partial y}
\left(\varrho^2 D^o\frac{\partial \psi_o}{\partial y}\right)
-
n
\frac{\Omega}{\varrho}
%\frac{\partial\rho D^m \tfrac{\partial\rho Y^m u_i}{\partial x_i}u_i}{\partial x_i}=
,
\end{equation} 
  
the fuel (F), 

\begin{equation} 
 \hat{U}\frac{\partial \psi_{F}}{\partial x}+
 \hat{V}\frac{\partial \psi_{F}}{\partial y}
=
\frac{1}{PrLe_F}\frac{\partial}{\partial y}
\left(
       \varrho^2 D^F\frac{\partial \psi_F}{\partial y}
\right)
- 
\frac{\Omega}{\varrho}
%\frac{\partial\rho D^m \tfrac{\partial\rho Y^m u_i}{\partial x_i}u_i}{\partial x_i}=
\end{equation} 

and the product of the reaction  

\begin{equation} 
 \hat{U}\frac{\partial \psi_{P}}{\partial x}+
 \hat{V}\frac{\partial \psi_{P}}{\partial y}
=
\frac{1}{PrLe_P}\frac{\partial}{\partial y}
\left(
       \varrho^2 D^P\frac{\partial \psi_P}{\partial y}
\right)
+ 
(1+n)
\frac{\Omega}{\varrho}
%\frac{\partial\rho D^m \tfrac{\partial\rho Y^m u_i}{\partial x_i}u_i}{\partial x_i}=
.
\end{equation} 

The momentum equation in the $x$ direction:
  
\begin{equation} 
 \hat{U}\frac{\partial \hat{U}}{\partial x}+
 \hat{V}\frac{\partial \hat{U}}{\partial y}=
\frac{\partial}{\partial y}\left(\mu\varrho\frac{\partial \hat{U}}{\partial  y}\right)
\end{equation} 
  
and energy equation 
  
\begin{equation}
\begin{split}
\hat{U}\frac{\partial }{\partial x}(c_p \overline{T})+
\hat{V}
\frac{\partial }{\partial y}(c_p \overline{T})=
\varrho
(\gamma-1)M^2\mu
\left(
\frac{\partial \hat{U}}{\partial y}
\right)
^2
+
\\
\frac{1}{Pr}\frac{\partial}{\partial y}\left(k\varrho\frac{\partial \overline{T}}{\partial y}\right)
-
\frac{Y^{i\infty}}{Pr}
\frac{\partial}{\partial y}
       \left(
              \frac{1}{Le_i}
              \sum_i\varrho^2 D_i\overline{h}_i\frac{\partial \psi_i}{\partial y}
       \right)
+
\frac{Q\Omega}{\varrho}
\end{split}
\end{equation} 

with $c_p=\sum_i \psi_i c_{p,i}$.

The energy equation can be rewritten in terms of the total enthalpy, 

\begin{equation}
H=\overline{h}+\frac{1}{2}\hat{U}^2
=c_p \overline{T}+\frac{1}{2}\hat{U}^2. 
\end{equation}

Multiplying the momentum equation by $\hat{U}$ 

\begin{equation} 
\hat{U}
 \hat{U}\frac{\partial \hat{U}}{\partial x}
 +
\hat{U}
 \hat{V}\frac{\partial \hat{U}}{\partial y}
 =
\hat{U}
\frac{\partial}{\partial y}\left(\mu\varrho\frac{\partial \hat{U}}{\partial  y}\right)
\end{equation} 

and adding to the energy
equation, we obtain

\begin{equation} 
\begin{split}
\hat{U}
 \hat{U}\frac{\partial \hat{U}}{\partial x}
 +
\hat{U}
 \hat{V}\frac{\partial \hat{U}}{\partial y}
+
\hat{U}\frac{\partial }{\partial x}(c_p \overline{T})
+
\hat{V}
\frac{\partial }{\partial y}(c_p \overline{T})
 =
 \\
\hat{U}
\frac{\partial}{\partial y}\left(\mu\varrho\frac{\partial \hat{U}}{\partial  y}\right)
+
\varrho
(\gamma-1)M^2\mu
\left(
\frac{\partial \hat{U}}{\partial y}
\right)
^2
+
\\
\frac{1}{Pr}\frac{\partial}{\partial y}\left(k\varrho\frac{\partial \overline{T}}{\partial y}\right)
-
\frac{Y^{i\infty}}{Pr}
\frac{\partial}{\partial y}
       \left(
              \frac{1}{Le_i}
              \sum_i\varrho^2 D_i\overline{h}_i\frac{\partial \psi_i}{\partial y}
       \right)
+
\frac{Q\Omega}{\varrho}
\end{split}
\end{equation}       

\begin{equation} 
\begin{split}
\hat{U}
 \hat{U}\frac{\partial \hat{U}}{\partial x}
 +
\hat{U}
 \hat{V}\frac{\partial \hat{U}}{\partial y}
+
\hat{U}\frac{\partial }{\partial x}(c_p \overline{T})
+
\hat{V}
\frac{\partial }{\partial y}(c_p \overline{T})
 =
 \\
\frac{\partial}{\partial y}\left(\mu \varrho\hat{U}\frac{\partial \hat{U}}{\partial  y}\right)
-
\mu 
\varrho
\left(
\frac{\partial \hat{U}}{\partial y}
\right)
^2
+
\varrho
(\gamma-1)M^2\mu
\left(
\frac{\partial \hat{U}}{\partial y}
\right)
^2
+
\\
\frac{1}{Pr}\frac{\partial}{\partial y}
\left(kc_p\varrho\frac{\partial \overline{h}}{\partial y}\right)
-
\frac{Y^{i\infty}}{Pr}
\frac{\partial}{\partial y}
       \left(
              \frac{1}{Le_i}
              \sum_i\varrho^2 D_i\overline{h}_i\frac{\partial \psi_i}{\partial y}
       \right)
+
\frac{Q\Omega}{\varrho}
\end{split}
\end{equation}       

Note that  

\begin{equation}
\overline{h}=\sum_{i} \psi_i \overline{h}_i
\end{equation}

then, 

\begin{equation}
\frac{\partial \overline{h}}{\partial y}
=
\sum_{i} \psi_i \frac{\partial \overline{h}_i}{\partial y}
+
\sum_{i} 
\frac{\partial \psi_i}{\partial y} \overline{h}_i
.
\end{equation}

Therefore,

\begin{equation} 
\begin{split}
(\rho u)\frac{\partial H}{\partial x}    
+
(\rho v)\frac{\partial H}{\partial y}    
=
 \\
\frac{\partial}{\partial y}
\left(
       \mu  \varrho\hat{U}\frac{\partial \hat{U}}{\partial  y}
 \right)
-
\mu 
\varrho
\left(
1-
(\gamma-1)M^2
\right)
\left(
\frac{\partial \hat{U}}{\partial y}
\right)
^2
\\
\frac{1}{Pr}
\frac{\partial}{\partial y}
\left[
       k\varrho c_p
       \left(
       \sum_{i} \psi_i \frac{\partial \overline{h}_i}{\partial y}
       +
       \sum_{i} 
       \frac{\partial \psi_i}{\partial y} \overline{h}_i
       \right)
\right]
-
\frac{Y^{i\infty}}{Pr}
\frac{\partial}{\partial y}
       \left(
              \frac{1}{Le_i}
              \sum_i\varrho^2 D_i\overline{h}_i\frac{\partial \psi_i}{\partial y}
       \right)
+
\frac{Q\Omega}{\varrho}
\end{split}
\end{equation}       

\begin{equation} 
\begin{split}
(\rho u)\frac{\partial H}{\partial x}    
+
(\rho v)\frac{\partial H}{\partial y}    
=
 \\
\frac{\partial}{\partial y}\left(\mu \varrho\hat{U}\frac{\partial \hat{U}}{\partial  y}\right)
-
\mu 
\varrho
\left(
1-
(\gamma-1)M^2
\right)
\left(
\frac{\partial \hat{U}}{\partial y}
\right)
^2
\\
\frac{1}{Pr}
\frac{\partial}{\partial y}
\left[
       k\varrho c_p
       \sum_{i} \psi_i \frac{\partial \overline{h}_i}{\partial y}
\right]
+
\frac{1}{Pr}
\frac{\partial}{\partial y}
\left[
       k\varrho c_p
       \sum_{i} 
       \frac{\partial \psi_i}{\partial y} \overline{h}_i
\right]
-
\frac{Y^{i\infty}}{Pr}
\frac{\partial}{\partial y}
\left(\sum_i\frac{1}{Le_i}\varrho^2 D_i\frac{\partial \psi_i}{\partial y}\overline{h}_i\right)
+
\frac{Q\Omega}{\varrho}
\end{split}
\end{equation}       

\begin{equation} 
\begin{split}
(\rho u)\frac{\partial H}{\partial x}    
+
(\rho v)\frac{\partial H}{\partial y}    
=
 \\
\frac{\partial}{\partial y}\left(\mu \varrho\hat{U}\frac{\partial \hat{U}}{\partial  y}\right)
-
\mu 
\varrho
\left(
1-
(\gamma-1)M^2
\right)
\left(
\frac{\partial \hat{U}}{\partial y}
\right)
^2
\\
\frac{1}{Pr}
\frac{\partial}{\partial y}
\left[
       k\varrho c_p
       \sum_{i} \psi_i \frac{\partial \overline{h}_i}{\partial y}
\right]
+
\frac{1}{Pr}
\left\{
\frac{\partial}{\partial y}
\left[
kc_p
\varrho
\left(
       \sum_{i} 
       \left(
       1
       -
       \frac{1}{Le_i}
       \frac{Y^{i\infty}
       \varrho D_i}{kc_p}
       \right)
       \frac{\partial \psi_i}{\partial y} \overline{h}_i
\right)
\right]
\right\}
+
\frac{Q\Omega}{\varrho}
\end{split}
\end{equation}       


\subsection{Chemical Reaction}

Considering a single step chemical reaction,

\begin{equation}
\nu_F F + \nu_O O \rightarrow \nu_P P
\end{equation}

where $\nu_F$ and $\nu_O$
are the stiochiometric coeficients, for the 
fuel and oxidizer, respectively. 
Considering $\nu_F=1$ and $\nu_O=n$, indicating that $n$ moles of oxidizer are
required to burn one mole of fuel, producing $\nu_p=(1+n)$ moles of product,

\begin{equation}
F + n O \rightarrow (1+n)P
\end{equation}

where $n$ is the stiochiometric 
mass ration between the fuel and oxidizer, which can be expressed as

\begin{equation}
n=\frac{W_F}{W_O}\frac{\nu_O}{\nu_F}
.
\end{equation}
Where $W_F$ and $W_O$ are the molecular weights of the fuel and oxidizer, respectively.

The reaction rate can be written as
\begin{equation}
\Omega=k\psi_F^a \psi_O^b
\end{equation}

where $a$ and $b$ are the order of the reaction with respect to the fuel and
oxidizer, respectively. $k$ is the reaction rate constant, which can be expressed
as a function of the temperature using the Arrhenius law:

\begin{equation}
k=A\exp\left(-\frac{E_a}{RT}\right)
\end{equation}

Zeldovich proposed using the $Z$ varible, which is a linear combination of the
mass fractions of the fuel and oxidizer, to reduce the number of
equations. Definning the passive scalars:

\begin{equation}
%Z=\frac{\psi_F-\psi_O+\psi_{O\infty}}{\psi_{F\infty}+\psi_{O\infty}}
Z_1=
\psi_F
-
\frac{\psi_O}{n}
\end{equation}          

which is a linear combination of the fuel and oxidizer mass fractions, 
meaning that for stiochiometric conditions, 
measures how far the mixture is from stoichiometric balance.
if $Z_1=0$ the mixture is stoichiometric, for a 
fueled mixture, $Z_1>1$, and for an oxidizer rich mixture, $Z_1<1$, 

Similarly, another passive scalar can be defined as
\begin{equation}
%Z=\frac{\psi_F-\psi_O+\psi_{O\infty}}{\psi_{F\infty}+\psi_{O\infty}}
Z_2=
\psi_F
+
\frac{\psi_P}{n+1}
\end{equation}          

which is a linear combination of the fuel and product mass fractions,
measuring how far the mixture is from the pure fuel. 
$Z_2=0$ only for a pure oxidizer flow. 
 
For the oxidizer, another passive scalar can be defined as

\begin{equation}
Z_3=
\psi_O
+
\frac{\psi_P}{n}
\end{equation}          

which is a linear combination of the oxidizer and product mass fractions,
measuring how far the mixture is from the pure oxidizer. 
$Z_3=0$ only for a pure fuel flow. 

Then, the species equation in term of passive scalars $z_i$ can be writen as  :  

\begin{equation} 
\begin{split} 
 \hat{U}\frac{\partial \psi_{F}}{\partial x}
 +
 \hat{V}\frac{\partial \psi_{F}}{\partial y}
 -
 \frac{1}{n}
 \hat{U}\frac{\partial \psi_{o}}{\partial x}
 -
 \frac{1}{n}
 \hat{V}\frac{\partial \psi_{o}}{\partial y}
=
\\
\frac{1}{PrLe_F}\frac{\partial}{\partial y}
\left(\varrho^2 D^F\frac{\partial \psi_F}{\partial y}\right)
+
\frac{s^{F}\Omega}{\varrho}
-
\frac{1}{n}
\frac{1}{PrLe_o}\frac{\partial}{\partial y}
\left(
       \varrho^2 D^o\frac{\partial \psi_o}{\partial y}
\right)
-
\frac{s^{o}\Omega}{\varrho}
+
,
\end{split} 
\end{equation} 
  
\begin{equation} 
\begin{split} 
 \hat{U}\frac{\partial Z_1}{\partial x}+
 \hat{V}\frac{\partial Z_1}{\partial y}
=
\\
\frac{1}{PrLe_o}\frac{\partial}{\partial y}
\left(
       \varrho^2
       D^o
       \left(
       \frac{D^F}{D^O}
       \frac{\partial \psi_F}{\partial y}
       -
       \frac{1}{n}
       \frac{\partial \psi_o}{\partial y}
       \right)
\right)
+
\frac{s^{F}\Omega}{\varrho}
-
\frac{1}{n}
\frac{s^{o}\Omega}{\varrho}
,
\end{split} 
\end{equation} 

Assuming that $D^F=D^O=D$ 
and $Le_O=Le_F=Le$

\begin{equation} 
\begin{split} 
 \hat{U}\frac{\partial Z_1}{\partial x}+
 \hat{V}\frac{\partial Z_1}{\partial y}
=
\frac{1}{PrLe}\frac{\partial}{\partial y}
\left(
       \varrho^2
       D
       \frac{\partial Z_1 }{\partial y}
\right)
+
\frac{\Omega}{\varrho}
-
\frac{1}{n}
n
\frac{\Omega}{\varrho}
,
\end{split} 
\end{equation} 

\begin{equation} 
\begin{split} 
 \hat{U}\frac{\partial Z_1}{\partial x}+
 \hat{V}\frac{\partial Z_1}{\partial y}
=
\frac{1}{PrLe}\frac{\partial}{\partial y}
\left(
       \varrho^2
       D
       \frac{\partial Z_1 }{\partial y}
\right)
,
\end{split} 
\end{equation} 

Similarly for $Z_2$:

\begin{equation} 
\begin{split} 
 \hat{U}\frac{\partial Z_2}{\partial x}+
 \hat{V}\frac{\partial Z_2}{\partial y}
=
\frac{1}{PrLe}\frac{\partial}{\partial y}
\left(
       \varrho^2
       D
       \frac{\partial Z_2 }{\partial y}
\right)
,
\end{split} 
\end{equation} 

with  $D^P=D$ 
and $Le_P=Le$. 
Finally with $Z_3$

\begin{equation} 
\begin{split} 
 \hat{U}\frac{\partial Z_3}{\partial x}+
 \hat{V}\frac{\partial Z_3}{\partial y}
=
\frac{1}{PrLe}\frac{\partial}{\partial y}
\left(
       \varrho^2
       D
       \frac{\partial Z_3 }{\partial y}
\right)
.
\end{split} 
\end{equation} 